\documentclass[10pt,twocolumn]{article}
\usepackage[T1]{fontenc}
\usepackage[utf8]{inputenc}
\usepackage{lmodern}
\usepackage[margin=1.8cm,top=2.4cm,bottom=2cm,columnsep=0.6cm]{geometry}
\usepackage{fancyhdr}
\usepackage{titlesec}
\usepackage{booktabs}
\usepackage{amsmath,amssymb,amsfonts}
\usepackage{microtype}
\usepackage{xcolor}
\usepackage{mdframed}
\usepackage{parskip}
\usepackage{enumitem}
\usepackage{hyperref}

\definecolor{rulered}{RGB}{180,30,30}
\definecolor{boxgray}{RGB}{245,245,245}
\definecolor{keywordgray}{RGB}{100,100,100}

\pagestyle{fancy}
\fancyhf{}
\fancyhead[L]{\textsc{\small Claude Mag}}
\fancyhead[C]{\small\textit{Machine Learning Research Digest}}
\fancyhead[R]{\small 2026-08-05}
\fancyfoot[C]{\small\thepage}
\renewcommand{\headrulewidth}{0.8pt}

\titleformat{\section}{\normalsize\bfseries\color{rulered}}{}{0pt}{}%
  [\vspace{-4pt}\color{rulered}\rule{\columnwidth}{0.4pt}]
\titlespacing{\section}{0pt}{8pt}{4pt}

\hypersetup{colorlinks=true,urlcolor=rulered,linkcolor=black}

\begin{document}
\setlength{\parindent}{0pt}
\setlength{\parskip}{4pt}

{\small\color{keywordgray}\textbf{Keywords:}
  multimodal reasoning \textbullet{}
  auxiliary memory \textbullet{}
  chain-of-thought \textbullet{}
  training-free inference}\par\vspace{4pt}

{\LARGE\bfseries\raggedright
  TRAM: Enhancing Multimodal Reasoning with Trajectory-Derived Auxiliary Memory}\par\vspace{4pt}

{\large\itshape\raggedright
  A Training-Free Inference Wrapper Recovers Reasoning-Derived Information Lost
  in Long Multimodal Chains Without Any Fine-Tuning}\par\vspace{6pt}

{\small\textbf{By} Kang Liu, Zijing Wang, Yongkang Liu, Mengjie Zhao,
  Xiaocui Yang, Shi Feng, Yifei Zhang \& Daling Wang}\par\vspace{2pt}

{\small\color{keywordgray}arXiv:2608.01922 \textbullet{} Submitted 3 August 2026}
\par\vspace{2pt}\noindent{\color{rulered}\rule{\columnwidth}{0.6pt}}\vspace{6pt}

As multimodal reasoning chains grow long, intermediate conclusions lose influence
over later steps. \textit{TRAM} fixes this without retraining: a training-free
inference wrapper extracts a compact trajectory memory from each intermediate step
and injects it into decoding, closing the accuracy gap caused by long-chain
information loss in Multimodal Large Reasoning Models (MLRMs).

\begin{mdframed}[backgroundcolor=boxgray,linecolor=rulered,linewidth=1pt,
    innerleftmargin=6pt,innerrightmargin=6pt,
    innertopmargin=6pt,innerbottommargin=6pt]
\textbf{\small AT A GLANCE}\par\vspace{4pt}
\begin{description}[leftmargin=0pt,labelsep=4pt,itemsep=2pt,parsep=0pt,
    font=\small\bfseries]
  \item[Problem:] In long reasoning chains, intermediate conclusions and
    visual inferences established early in the trajectory lose influence over
    later decoding steps, causing accuracy to degrade.
  \item[Why it matters:] Hard visual reasoning benchmarks require chains of
    many steps; state-of-the-art MLRMs degrade precisely on these long-chain
    cases where the capability would matter most.
  \item[Method:] Training-free auxiliary memory extracted from the reasoning
    trajectory so far and injected as additional context at each decoding step.
  \item[Result:] Improved accuracy on long-chain multimodal reasoning benchmarks
    without fine-tuning or model weight modification.
  \item[Evidence:] Attribution analysis showing correctness correlates with
    trajectory retention, not visual attribution alone; benchmark improvements
    across multiple MLRMs.
\end{description}
\end{mdframed}

\section{The Big Picture}

Multimodal Large Reasoning Models (MLRMs) combine a vision encoder with a large
language model to tackle tasks that require reading an image, constructing
multi-step inferences, and arriving at a final answer. Representative models
include LLaVA-o1, InternVL-Reasoner, and similar 2026 state-of-the-art systems.

These models work well on short reasoning tasks but degrade on long chains---precisely
the cases where step-by-step reasoning is most valuable. TRAM provides a training-free
remedy, applicable as a plug-in wrapper around any deployed MLRM.

\section{Why Existing Approaches Fall Short}

\textbf{Retraining with longer context windows} is expensive and risks catastrophic
forgetting of other capabilities. Models with 32K+ context windows still degrade
on long reasoning tasks because the issue is not available context length---the
intermediate conclusions \emph{are} in the context---but effective attention to them.

\textbf{Fine-tuning on curated long-chain examples} is data-hungry, hard to generalise,
and requires modifying model weights. This prevents deployment on proprietary models
accessed via API.

\textbf{Prompt engineering} (e.g., adding explicit ``remember that...'' prompts) is
fragile and task-specific.

TRAM avoids all of these by operating as an inference-time wrapper with no access
to model weights.

\section{Core Mechanism}

TRAM rests on an attribution insight: correctness in MLRMs is not well-predicted by
visual token attribution (how much the model ``looks at'' the image), but by whether
the trajectory up to each step retains and integrates earlier reasoning-derived
information. This means the bottleneck is not the visual encoder---it is the
propagation of intermediate conclusions through the reasoning chain.

TRAM addresses this by maintaining an explicit \textbf{trajectory memory} $\mathcal{M}$
that accumulates reasoning-derived information as generation proceeds:

\begin{enumerate}[itemsep=2pt,parsep=0pt]
  \item At each completed reasoning step $s_t$, a lightweight extractor parses the
    step and identifies key propositions (visual relations, numeric facts, logical
    constraints, sub-conclusions).
  \item These are added to $\mathcal{M}$ indexed by type and step index.
  \item When generating the next step $s_{t+1}$, the most relevant entries from
    $\mathcal{M}$ are retrieved and prepended to the decoding context as an
    ``Auxiliary Memory'' block.
\end{enumerate}

The extractor and retriever are lightweight and rule-based; no neural modules
beyond the MLRM itself are required.

\section{Technical Formulation}

Let $\mathcal{H}_t = (v, s_1, s_2, \ldots, s_t)$ be the full trajectory at step~$t$,
where $v$ is the visual input and each $s_i$ is a reasoning step. The MLRM
without TRAM generates:
\[
s_{t+1} \sim p_\theta(s_{t+1} \mid \mathcal{H}_t).
\]
TRAM augments this with trajectory memory $\mathcal{M}_t$, a structured
summary of $\mathcal{H}_t$:
\[
s_{t+1} \sim p_\theta\!\left(s_{t+1} \mid \mathcal{H}_t, \mathcal{M}_t\right),
\]
where $\mathcal{M}_t$ is constructed by:
\[
\mathcal{M}_t = \bigoplus_{i=1}^{t} \mathrm{Extract}(s_i),
\]
with $\mathrm{Extract}$ parsing the step text and $\bigoplus$ denoting
structured accumulation into the memory store.

The attribution score used to motivate the design is:
\[
A_{\mathrm{traj}}(s_j, s_{t+1}) = \sum_{k} \alpha_{k,j}^{(t)},
\]
where $\alpha_{k,j}^{(t)}$ is the attention weight from position $k$ in
step $s_{t+1}$ to token $j$ in a prior step. The empirical finding is that
$A_{\mathrm{traj}}$ for early steps decreases as $t$ grows, explaining the
long-chain accuracy degradation.

\section{Learning or Inference Procedure}

TRAM operates as a post-processing wrapper around any MLRM's generation loop:
\begin{enumerate}[itemsep=2pt,parsep=0pt]
  \item Load the MLRM (weights unmodified).
  \item Initialise $\mathcal{M}_0 = \emptyset$.
  \item For each reasoning step $t$:
    \begin{enumerate}[itemsep=1pt,parsep=0pt]
      \item Retrieve top-$K$ entries from $\mathcal{M}_{t-1}$ relevant to the
        current partial step (keyword + type matching).
      \item Prepend retrieved entries as an ``Auxiliary Memory'' block to the
        MLRM's context.
      \item Sample $s_t$ from the augmented context.
      \item Run $\mathrm{Extract}(s_t)$ and add results to $\mathcal{M}_t$.
    \end{enumerate}
  \item Return the final answer from the last reasoning step.
\end{enumerate}
No gradient computation or parameter update occurs at any step.

\section{What the Guarantee Says}

The method provides no formal convergence guarantee, but the attribution
analysis provides empirical evidence for the causal chain: (i) early trajectory
attention decays with chain length; (ii) TRAM injects a compressed representation
of the decayed information; (iii) accuracy on long-chain benchmarks improves.

The compression step (Extract) is necessarily lossy, so TRAM provides benefit
primarily when the extracted propositions are those the model would have failed
to recover from the diluted attention---a condition satisfied empirically for the
benchmark tasks studied.

\section{Experimental Findings}

Evaluated across multiple MLRMs on long-chain multimodal reasoning benchmarks:
\begin{itemize}[itemsep=2pt,parsep=0pt]
  \item Consistent accuracy improvement on long-chain reasoning tasks (those
    requiring $>6$ reasoning steps) across all tested MLRMs.
  \item Minimal or no improvement on short-chain tasks, consistent with the
    theory that attention dilution is not a bottleneck when chains are short.
  \item Improvement is larger for tasks requiring explicit recall of visual
    relations or numeric facts established early in the chain.
  \item Wall-clock overhead: the extractor and memory retrieval add negligible
    latency relative to the MLRM's own generation time.
\end{itemize}

\section{Ablations and Interpretation}

\textbf{Memory type:} All three memory entry types (visual relations, numeric
facts, logical constraints) contribute; removing any one reduces gains by 15--30\%.

\textbf{Retrieval top-$K$:} Performance peaks at $K=3$--5; beyond $K=8$ the
auxiliary block becomes too long and begins to hurt attention focus.

\textbf{Extractor quality:} Replacing the rule-based extractor with an LLM-based
extractor improves results slightly but increases latency. The rule-based version
provides most of the gain at negligible cost.

\textbf{Baseline---full context:} Simply prepending the full $\mathcal{H}_t$
history twice does not help; the structured compression step is essential.

\section*{Reference}

Kang Liu, Zijing Wang, Yongkang Liu, Mengjie Zhao, Xiaocui Yang,
Shi Feng, Yifei Zhang, Daling Wang.
\textbf{TRAM: Enhancing Multimodal Reasoning with Trajectory-Derived Auxiliary Memory}.
arXiv:2608.01922, August 2026.
\url{https://arxiv.org/abs/2608.01922}

\end{document}
